%=====================================================================
\chapter{GPU and Accelerator Virtualization}\label{ch:gpu}
%=====================================================================
\locator{6}{Part VI --- Emerging Hardware and Practice}

\begin{outcomes}
By the end of this chapter you will be able to:
\begin{tight}
  \item \textbf{Explain} why a GPU is harder to share than a processor, and
        \textbf{name} the three sharing models.
  \item \textbf{Distinguish} time-sliced vGPU from MIG partitioning by what is
        divided and what is guaranteed.
  \item \textbf{Compute} the utilisation and cost consequences of each model
        for a stated set of jobs.
  \item \textbf{Describe} what a DPU moves off the host, and \textbf{relate} it
        to \chref{paravirt}.
  \item \textbf{Choose} a sharing model for a workload and \textbf{state} what
        it gives up.
\end{tight}
\end{outcomes}

\begin{prereq}
\chref{paravirt} entire --- passthrough, VFIO, SR-IOV and the IOMMU are the
machinery this chapter applies to a different device --- and \chref{cpumem}
(scheduling and overcommit, which work differently here).
\end{prereq}

\begin{keyterms}
GPU passthrough $\bullet$ vGPU $\bullet$ time slicing $\bullet$ MIG $\bullet$
instance profile $\bullet$ SR-IOV for GPUs $\bullet$ framebuffer $\bullet$
context switch $\bullet$ mediated device $\bullet$ VFIO-mdev $\bullet$ DPU
$\bullet$ SmartNIC $\bullet$ infrastructure offload $\bullet$ HBM $\bullet$
NVLink $\bullet$ fractional GPU
\end{keyterms}

\section{Why a GPU Is Not a CPU}

\chref{cpumem} divided a processor among guests by scheduling: a vCPU runs for
a slice, is descheduled, and another runs. The state to save is a few hundred
bytes of registers, and the switch costs microseconds.

A GPU resists this for three reasons.

\begin{tight}
  \item \textbf{The state is enormous.} A context includes tens of gigabytes of
        high-bandwidth memory holding model weights, textures and intermediate
        results. You cannot save and restore that between time slices; there is
        nowhere to put it and no time to move it.
  \item \textbf{The work is long and uninterruptible.} A compute kernel runs to
        completion. Pre-empting it mid-flight is either impossible or very
        expensive, depending on the generation.
  \item \textbf{The memory is the scarce resource, not the compute.} A training
        job that needs 40 GB of GPU memory cannot run in 20 GB at half speed.
        It cannot run at all.
\end{tight}

That last point is the one that matters most. Processor overcommit works
because time is divisible (\chref{cpumem}); GPU memory is as indivisible as
host memory, and there is far less of it.

\dsfig{%
\begin{tikzpicture}[font=\scriptsize,
  gpu/.style={rounded corners=3pt, line width=1.1pt, minimum width=4.6cm,
              minimum height=2.6cm},
  seg/.style={rounded corners=1.5pt, line width=0.6pt, align=center,
              font=\tiny, inner sep=1pt},
  hd/.style={font=\scriptsize\bfseries, align=center},
  note/.style={font=\scriptsize\itshape, align=center}
]
  % --- passthrough ---
  \node[gpu, draw=purplehl, fill=purplehl!5] at (0,0) {};
  \node[hd, text=purplehl!70!black] at (0,1.55) {Passthrough};
  \node[seg, draw=purplehl, fill=purplehl!22, text=purplehl!70!black,
        minimum width=3.9cm, minimum height=1.7cm] at (0,-0.1)
       {one guest\\[2pt] 80 GB \ $\bullet$\ all SMs\\[2pt]
        \textbf{native speed}};
  \node[note, text=purplehl!70!black] at (0,-1.75)
       {1 guest per card};

  % --- time sliced vGPU ---
  \node[gpu, draw=goldhl, fill=goldhl!5] at (5.6,0) {};
  \node[hd, text=goldhl!45!black] at (5.6,1.55) {vGPU (time sliced)};
  \foreach \i in {0,1,2,3}{
    \pgfmathsetmacro{\y}{0.62-\i*0.44}
    \node[seg, draw=goldhl, fill=goldhl!18, text=goldhl!45!black,
          minimum width=3.9cm, minimum height=0.38cm] at (5.6,\y)
         {guest \number\numexpr\i+1\relax \ --- 20 GB, \emph{all} SMs in turn};}
  \node[note, text=goldhl!45!black] at (5.6,-1.75)
       {memory split, compute shared\\in time: no guarantee};

  % --- MIG ---
  \node[gpu, draw=greenhl, fill=greenhl!5] at (11.2,0) {};
  \node[hd, text=greenhl!45!black] at (11.2,1.55) {MIG (partitioned)};
  \foreach \i in {0,1,2,3}{
    \pgfmathsetmacro{\y}{0.62-\i*0.44}
    \node[seg, draw=greenhl, fill=greenhl!18, text=greenhl!45!black,
          minimum width=3.9cm, minimum height=0.38cm] at (11.2,\y)
         {guest \number\numexpr\i+1\relax \ --- 20 GB + \emph{its own} SMs};}
  \node[note, text=greenhl!45!black] at (11.2,-1.75)
       {memory \emph{and} compute split\\in hardware: guaranteed};

  \node[rounded corners=3pt, draw=accent!55, fill=accent!5,
        text=accent!70!black, font=\scriptsize, align=center, inner sep=4pt,
        text width=12.6cm] at (5.6,-2.75)
       {The difference between the middle and the right is whether a noisy
        neighbour exists. Under time slicing one guest's long kernel delays
        the others; under MIG it cannot, because the streaming
        multiprocessors are not shared.};
\end{tikzpicture}}%
{Three ways to divide one accelerator. Only the right-hand one gives a
guarantee, and only the left-hand one gives the whole card.}{gpu-sharing}

\section{The Three Models}

\begin{definitionbox}[Passthrough, vGPU and MIG]
\term{Passthrough} assigns the whole physical GPU to one guest with VFIO and
the IOMMU, exactly as \chref{paravirt} assigned a network card. Native
performance, the guest's own vendor driver, and all of \chref{paravirt}'s
costs: no live migration, memory pinned, one guest per card.

\smallskip
\term{vGPU} --- NVIDIA's term; a \term{mediated device} in Linux terms --- splits
the card's memory between several guests and \emph{time slices} the compute
between them. Each guest gets a fixed framebuffer and a turn at all the
streaming multiprocessors. Density without a performance guarantee.

\smallskip
\term{MIG} --- multi-instance GPU, from the A100 onward --- partitions the
hardware itself. Memory, cache, streaming multiprocessors and memory bandwidth
are divided into up to seven instances with hardware-enforced boundaries. Each
instance behaves like a smaller, independent GPU. Density \emph{with} a
guarantee, at the cost of fixed partition sizes.
\end{definitionbox}

\begin{conceptbox}[Which One Matches Which Workload]
The deciding question is whether the workload's demand is steady or bursty, and
whether it can tolerate interference.

\begin{tight}
  \item \textbf{Training a large model} --- passthrough, or MIG's largest
        profile. It needs all the memory and it runs for days; sharing buys
        nothing.
  \item \textbf{Inference serving} --- MIG. Many small, latency-sensitive
        requests that must not be delayed by somebody else's batch job. The
        guarantee is the point.
  \item \textbf{Virtual desktops with CAD or rendering} (\chref{desktop}) ---
        vGPU. Demand is bursty and interactive, users are rarely busy
        simultaneously, and time slicing exploits exactly that.
  \item \textbf{Student or research workloads} --- MIG where the hardware
        supports it, because fairness between users matters more than peak
        throughput, and one runaway job must not stop the others.
\end{tight}
\end{conceptbox}

\begin{worked}[Four Jobs, One Card]
A department has one 80 GB accelerator and four research groups. Each needs
18 GB of GPU memory and runs jobs averaging 40 minutes, roughly 6 hours a day
each. Compare passthrough, vGPU and MIG.

\smallskip
\textbf{Step 1 --- passthrough.} One guest owns the card.
\[ \text{concurrent groups} = 1 \qquad
   \text{utilisation} = \frac{6}{24} = \mathbf{25\%} \]
The other three queue. Aggregate demand is $4 \times 6 = 24$ hours a day
against 24 hours available, so the card is theoretically just sufficient ---
but only if the four groups never want it at the same time, which they will,
because they work the same hours. Expect long queues at 2 p.m. and an idle card
at 3 a.m.

\smallskip
\textbf{Step 2 --- vGPU, four slices of 20 GB.} All four run concurrently.
Memory fits ($4 \times 20 = 80$). Compute is time sliced, so when all four are
busy each gets roughly a quarter of the throughput:
\[ \text{job time when all four are active} = 40 \times 4
   = \mathbf{160 \text{ minutes}} \]
When only one is active it gets the whole card and finishes in 40. Throughput
is good; \emph{predictability is gone}, and a researcher cannot tell whether
their job will take 40 minutes or 160.

\smallskip
\textbf{Step 3 --- MIG, four instances.} On an 80 GB card the
\texttt{2g.20gb} profile gives four instances of 20 GB with a fixed share of
compute --- about a quarter each, always.
\[ \text{job time, always} = \frac{40}{0.25} = \mathbf{160 \text{ minutes}} \]
The same 160 minutes as the worst vGPU case --- but it is \emph{always} 160,
whether the neighbours are busy or idle.

\smallskip
\textbf{Step 4 --- compare honestly.}
\begin{tight}
  \item Passthrough: fastest single job (40 min), 25\% utilisation, queueing.
  \item vGPU: 40--160 min depending on neighbours, high utilisation,
        unpredictable.
  \item MIG: always 160 min, high utilisation, predictable, and no job can be
        harmed by another.
\end{tight}

\smallskip
\textbf{Step 5 --- the recommendation, and why it is not obvious.} For four
research groups sharing one departmental card, \textbf{MIG}. The reason is not
throughput --- vGPU matches or beats it --- but that MIG makes the system
\emph{explicable}. A researcher who is told ``you have a quarter of a card and
your job takes 160 minutes'' can plan. One whose job takes 40 minutes on Tuesday
and 160 on Wednesday for reasons they cannot see will conclude the cluster is
broken, and will start asking for their own card.

\smallskip
\textbf{Step 6 --- and when to reverse it.} If the four groups' work were
genuinely bursty and rarely simultaneous --- interactive visualisation rather
than batch training --- vGPU would win outright, because time slicing harvests
exactly the idleness that MIG's fixed partitions would waste.
\end{worked}

\begin{alertbox}[Three practical constraints that decide this before you do]
\textbf{Licensing.} NVIDIA vGPU requires a separate software licence per
concurrent user, renewed annually. It is frequently a larger cost than the
card over three years, and it is the reason many sites use passthrough or MIG
instead.

\smallskip
\textbf{MIG profiles are fixed.} You choose from a published list --- one
instance of seven-eighths, two of a half, seven of a seventh, and a few
combinations. You cannot ask for 30\% of a card. Reconfiguring requires
stopping every workload on it.

\smallskip
\textbf{Not every card supports every model.} MIG arrived with the A100 and is
absent from consumer cards entirely; vGPU needs a supported card and the
licence; passthrough works almost everywhere. Check the specific model before
designing around a capability.
\end{alertbox}

\section{DPUs: Moving the Infrastructure Off the Host}

The last few chapters have quietly accumulated work the host performs on behalf
of guests: the virtual switch of \chref{network}, the encapsulation of VXLAN,
the storage path of \chref{storage}, the encryption everywhere. On a busy host
this is a material fraction of the processor --- capacity bought for guests and
spent on plumbing.

A \term{DPU} (data processing unit, or SmartNIC) is a card with its own
processors, memory and operating system that takes that work. The virtual
switch runs on the card. The storage initiator runs on the card. Encryption
runs on the card. The host's cores go back to the guests.

\begin{conceptbox}[Why This Is a Security Change as Much as a Performance One]
The obvious benefit is reclaimed cores --- commonly 20--30\% of a host that was
doing heavy overlay networking.

\smallskip
The structural benefit is that the infrastructure now runs on a processor the
host's operating system does not control. A compromised hypervisor cannot
rewrite the virtual switch's rules if the switch is not on the hypervisor; the
management path to the DPU is separate. This is \chref{platforms}'s
disaggregation argument --- Xen's driver domains --- implemented in silicon
rather than in software, and it is why AWS's Nitro architecture is built this
way.

\smallskip
It also completes an arc this book has been tracing since \chref{paravirt}:
emulate the device, then paravirtualize it, then hand the real one over, and
finally move the whole supporting apparatus to a computer of its own.
\end{conceptbox}

\begin{pitfall}
\begin{tight}
  \item \textbf{Assuming GPU memory can be overcommitted.} It cannot. A job
        needing 40 GB in a 20 GB instance does not run slowly; it fails.
  \item \textbf{Choosing vGPU for latency-sensitive inference.} Time slicing
        has no guarantee, and the tail latency of \chref{performance} is
        exactly what suffers.
  \item \textbf{Forgetting the vGPU licence in the budget.} It can exceed the
        hardware cost over three years.
  \item \textbf{Designing for MIG on a card that does not have it.} Check the
        model, not the vendor.
  \item \textbf{Expecting to live-migrate a GPU guest.} All of
        \chref{paravirt}'s passthrough costs apply, and GPU state makes them
        worse.
  \item \textbf{Measuring GPU utilisation as a percentage and stopping there.}
        ``100\% utilisation'' from \texttt{nvidia-smi} means a kernel was
        resident, not that the hardware was fully used. Look at achieved
        occupancy and memory bandwidth.
  \item \textbf{Buying accelerators before measuring queueing.} The worked
        example's passthrough case was 25\% utilised and still had queues,
        which is a scheduling problem, not a capacity one.
\end{tight}
\end{pitfall}

\begin{lab}[Partitioning an Accelerator]
Runs against a real GPU where one is available, and against captured output
where one is not --- the file is in \texttt{code/} and the exercise is the
same. Expect thirty minutes.

\begin{lstlisting}[language=bash]
#!/usr/bin/env bash
# ch20_gpu.sh -- passthrough, MIG and what utilisation really reports.
set -euo pipefail

# --- 1. Is there an accelerator here at all? ---------------------------
if ! command -v nvidia-smi >/dev/null; then
  echo "no GPU on this host -- read code/ch20_capture.txt instead"
  sed -n '1,40p' "$(dirname "$0")/ch20_capture.txt" 2>/dev/null || true
  exit 0
fi
nvidia-smi --query-gpu=name,memory.total,mig.mode.current \
  --format=csv

# --- 2. The IOMMU group: what passthrough would have to take. ---------
# Chapter 4: the whole group moves together, or nothing does.
for d in $(lspci -D | awk '/NVIDIA/{print $1}'); do
  grp=$(basename "$(dirname "$(readlink -f \
        "/sys/bus/pci/devices/$d/iommu_group")")")
  echo "device $d is in IOMMU group $grp"
  ls "/sys/kernel/iommu_groups/$grp/devices/"
done
# A card and its audio function usually share a group: both must go.

# --- 3. Enable MIG, if this card supports it. -------------------------
sudo nvidia-smi -i 0 -mig 1 || \
  echo "MIG not supported on this card -- see the alertbox"
sudo nvidia-smi mig -i 0 -lgip      # the fixed profiles available

# --- 4. Create four instances, and see them as separate devices. ------
sudo nvidia-smi mig -i 0 -cgi 2g.20gb,2g.20gb,2g.20gb,2g.20gb -C || true
nvidia-smi -L                        # four MIG UUIDs, one card

# --- 5. Run the same job on one instance, then on all four. -----------
JOB='python3 -c "
import time
try:
    import torch
except ImportError:
    raise SystemExit(\"install torch to run this step\")
a = torch.randn(8192, 8192, device=\"cuda\")
t = time.time()
for _ in range(50): a = a @ a.T / 1e4
torch.cuda.synchronize(); print(f\"{time.time()-t:.1f} s\")"'

mig_uuids () {
  nvidia-smi -L | awk -F'UUID: ' '/MIG/{print $2}' | tr -d ')'
}
UUID=$(mig_uuids | head -1)
echo "=== alone on one instance ==="
CUDA_VISIBLE_DEVICES="$UUID" bash -c "$JOB"

echo "=== all four instances busy together ==="
mig_uuids | while read -r u; do
  CUDA_VISIBLE_DEVICES="$u" bash -c "$JOB" &
done
wait
# Under MIG the single-instance time and the all-busy time should be
# the SAME. That is the guarantee the worked example is about.

# --- 6. What "utilisation" actually reports. --------------------------
nvidia-smi --query-gpu=utilization.gpu,utilization.memory,\
memory.used,memory.total --format=csv -l 1 -c 5
# utilization.gpu is "a kernel was resident during the sample", not
# "the hardware was fully used". See the pitfalls.

# --- 7. Tear down the partitions. -------------------------------------
sudo nvidia-smi mig -i 0 -dci || true
sudo nvidia-smi mig -i 0 -dgi || true
sudo nvidia-smi -i 0 -mig 0  || true
\end{lstlisting}

\textbf{What to notice.} Step 5 is the whole chapter in one measurement. Under
MIG the two timings match, because the partitions are enforced in hardware.
Repeat the same test under vGPU time slicing and they will not --- and that
difference, not throughput, is what you are choosing between.
\end{lab}

\begin{checkpoint}
\begin{tightnum}
  \item Explain why GPU memory cannot be overcommitted the way host memory can,
        referring to \chref{cpumem}'s reclamation ladder.
  \item Six groups each need 12 GB and run 5 hours a day on one 80 GB card.
        Compare passthrough and a six-way MIG partition on utilisation and on
        predictability.
  \item \texttt{nvidia-smi} reports 100\% GPU utilisation for a job that a
        colleague says is poorly optimised. Explain how both can be true, and
        what you would measure instead.
\end{tightnum}
\end{checkpoint}

\begin{chaptersummary}
\begin{tight}
  \item A GPU resists time slicing because its context is tens of gigabytes,
        its kernels are uninterruptible, and memory --- not compute --- is the
        scarce, indivisible resource.
  \item Passthrough gives one guest the whole card at native speed, with all of
        \chref{paravirt}'s costs.
  \item vGPU splits memory and time slices compute: density without a
        guarantee. MIG partitions memory \emph{and} compute in hardware:
        density with a guarantee, in fixed profiles.
  \item In the worked example MIG and vGPU gave the same throughput and only
        MIG gave a predictable answer --- which is usually what shared
        research and inference infrastructure needs.
  \item Licensing, fixed profiles and card support frequently decide the
        question before the engineering does.
  \item A DPU moves the virtual switch, the storage path and encryption to a
        computer of its own: 20--30\% of host cores returned, and the
        infrastructure placed outside the hypervisor's control.
\end{tight}
\end{chaptersummary}

\begin{reviewq}
  \item \textbf{(MCQ)} MIG differs from time-sliced vGPU in that it:
        (a)~supports more guests; (b)~partitions compute and memory in
        hardware, so performance is guaranteed; (c)~requires no licence;
        (d)~allows live migration.
  \item \textbf{(MCQ)} A job needing 40 GB of GPU memory, given a 20 GB
        instance, will: (a)~run at half speed; (b)~swap to host memory;
        (c)~fail; (d)~be queued automatically.
  \item \textbf{(Short)} Give the three reasons a GPU is harder to share than a
        CPU, and say which is the most consequential.
  \item \textbf{(Short)} State a workload suited to vGPU and one suited to MIG,
        with the reason for each.
  \item \textbf{(Short)} Explain what a DPU moves off the host and why that is
        a security change as well as a performance one.
  \item \textbf{(Applied)} A department will buy one accelerator for eight
        research students whose jobs need 9 GB each and run in bursts during
        office hours. Recommend a sharing model, give the expected utilisation,
        and state the two questions you would answer before purchase.
  \item \textbf{(Applied)} An inference service has a p99 latency requirement
        of 150 ms and currently runs on time-sliced vGPU, missing the target
        intermittently. Explain the mechanism, propose a change, and say what
        you would measure to confirm the fix.
\end{reviewq}
